\cotefile{148}
\scanpage{3}

Une façon de le voir est de dire que
la donnée de $\uline{s}$ définit un scindage
\inserted{$\varphi_{\uline{s}}$} de
l'extension $(*)$ ci-dessus, \unsure{et que vu pour|et en tant que}
ss-groupe de $\operatorname{SRT}(X)$,
$S_1T(X,\uline{s})$ s'identifie
à l'image \unsure{inverse|\deleted{d'}inverse} de
\[
\varphi_{\uline{s}}(T(X))\subset\operatorname{SUT}(X)
\]
par l'hom. can
$\operatorname{SUR}(X)\mathrel{\unsure{\text{---}|\longrightarrow}}\operatorname{SUT}(X)$
de passage au quotient. Le chgt de
$\uline{s}$ en $\uline{s}'=u.\uline{s}$
transforme $\varphi_{\uline{s}}$ en
\[
\varphi_{\uline{s}'}=\operatorname{int}(u)\circ\varphi_{\uline{s}}
\]
donc (si $\widetilde u\in\mathbb R^I$ relève $u$\unsure{,|)} on trouve
\[
S_1T(X,\uline{s}')
=\operatorname{int}(\widetilde u)S_1T(X,\uline{s})\unsure{.|,}
\]

En termes de $\unsure{\overline X|\vec X}$, la définition de
$S_1T(X,\uline{s})$
dépend un choix $\uline{s}$ \uline{non-anodin},
par contre celle de $S_1RT(X)$ ne dépend que d'
un choix anodin, celui de la famille des
structures toriques $(t_i)_{i\in I}$.
De plus $\prod T_i$ comme
torseur sous $\unsure{\mathbb U^I|\mathbb M^I}$ se réalise
\inserted{(pas fidèlement d'ailleurs)}
comme \unsure{une classe|un des} de scindages de l'ext. $(*)$
(via autom. int. par $\mathbb U^I$)
$\bigl[$déduits de $\operatorname{SRT}(X)$
avec les structures suppl.$\unsure{\bigr\}|\bigr]}$
savoir\unsure{)|,} les inclusions

\diagram{
\begin{tikzpicture}[baseline=(current bounding box.center)]
\node (a) at (0,1.65) {$\mathbb{R}^{I}$};
\node (b) at (0,-1.65) {$S_1T^{!}(X)$};
\node (c) at (4.6,0) {$S\unsure{R|K}T(X)$};
\draw[->] (a.east) to[out=-35,in=165] (c.north west);
\draw (3.12,.73) .. controls (3.35,.66) and (3.52,.43) .. (3.73,.24);
\draw[->] (b.north east)--(c.south west);
\end{tikzpicture}
}

\unsure{coïncid.|coïncidant|coincident} sur $\mathbb Z^I$.

\margin{2}
